% Transcribed human main text; editorial scope in tex/009.
\paragraph{Crossing number estimates}
\begin{proposition}[1.2]
The crossing number of $G$ is at least $E-3V$.
\end{proposition}
\begin{proof}
Let $k(G)$ be the crossing number of $G$. Embed $G$ in the plane with $k(G)$ crossings. By removing at most $k(G)$ edges, we get a planar graph $G'$ with $E'=E-k$ edges and $V'\leq V$ vertices. We see $0\geq E'-3V'\geq E-k-3V$.
\end{proof}

\begin{theorem}[1.3]
If $G$ is a graph with $E$ edges and $V$ vertices, and $E\geq4V$, then the crossing number of $G$ is at least $(1/64)E^3V^{-2}$.
\end{theorem}
This theorem was proven by several authors before Sz\'ekely, but we give his proof.
\begin{proof}
Let $p$ be a number between $0$ and $1$ which we choose below. Let $G'$ be a random subgraph of $G$ formed by including each vertex of $G$ independently with probability $p$. We include an edge of $G$ in $G'$ if its endpoints are in $G'$.

We consider the expected values for the number of vertices and edges in $G'$. The expected value of $V'$ is $pV$. The expected value of $E'$ is $p^2E$. For every subgraph $G'\subset G$, the crossing number of $G'$ is at least $E'-3V'$. Therefore, the expected value of the crossing number of $G'$ is at least $p^2E-3pV$.

On the other hand, we give an upper bound on the expected crossing number of $G'$ as follows. Let $k=k(G)$ be the crossing number of $G$. Let $F:G\to\mathbb R^2$ be a legal embedding with $k$ crossings. We claim that each crossing of $F$ involves two disjoint edges. In other words, two edges that share a vertex don't cross. We come back to the claim at the end. By restricting $F$ to $G'$, we get an embedding of $G'$ with $p^4k$ crossings on average. This is because each crossing involves four vertices, and it appears as a crossing of $F(G')$ only if all four vertices are included in $G'$. (If $F$ had a crossing involving two edges containing a common vertex, then it would appear with the much higher probability $p^3$.) Therefore, the expected value of the crossing number of $G'$ is at most $p^4k$.

Comparing our upper and lower bounds, we see that $p^4k\geq p^2E-3pV$, and so we get the following lower bound for $k$.
\[
k\geq p^{-2}E-3p^{-3}V.
\]
We can now choose $p$ to optimize the right-hand side. We choose $p=4V/E$, and we have $p\leq1$ since we assumed $4V\leq E$. Plugging in we get $k\geq(1/64)E^3V^{-2}$.

To finish the proof, we just have to check the claim that $F$ has no crossings of edges that share a vertex. Given any map with such a crossing, we explain how to modify it to reduce the crossing number. Say that $F(e_1)$ and $F(e_2)$ each leave $F(v)$ and cross at $x$. (If they cross several times, then let $x$ be the last crossing.) We modify $F$ as follows. Suppose that $F(e_1)$ crosses $k_1$ other edges on the way from $F(v)$ to $x$ and that $F(e_2)$ crosses $k_2$ other edges on the way from $F(v)$ to $x$. We choose the labelling so that $k_1\leq k_2$. Then we modify $F$ on the edge $e_2$, making $F(e_2)$ follow parallel to $F(e_1)$ until $x$ and then rejoin its original course at $x$, so that $F(e_1)$ and $F(e_2)$ never cross. This operation reduces the crossing number of $x$, and so a minimal map $F$ has no such crossings.
\end{proof}

\paragraph{The Szemer\'edi--Trotter theorem}
\begin{theorem}[2.1]
Let $\mathcal L$ be a set of $L$ lines in the plane. Let $P_k$ be the set of points that lie on at least $k$ lines of $\mathcal L$. Then the number of points in $P_k$ is at most
\[
\max(2Lk^{-1},2^9L^2k^{-3}).
\]
\end{theorem}
\begin{proof}
Using the lines and points, we make a graph mapped into the plane. The vertices of our graph $G$ are the points of $P_k$. We join two vertices with an edge of $G$ if the two points are two consecutive points of $P_k$ on a line $l\in\mathcal L$. This graph is not embedded, but the crossing number of our map is at most $\binom L2\leq L^2$, since each crossing of the graph $G$ must correspond to an intersection of two lines of $\mathcal L$.

We will count the vertices and edges of the graph $G$ and apply the crossing number theorem. The number of vertices of our graph is $V=|P_k|$. The number of edges of our graph is $kV-L$. (At first sight, each vertex should be adjacent to $2k$ edges which would give $kV$ edges. But on each line $l\in\mathcal L$, the first and last vertices are adjacent to one less edge than this initial count.) As long as $E\geq4V$, we can apply the crossing number theorem and it gives
\[
L^2\geq(1/64)(kV-L)^3V^{-2}.
\]
Either $V\leq2L/k$, or else $kV-L\geq(1/2)kV$. In the former case, we are done. In the latter case, we have $L^2\geq2^{-9}k^3V$, which means $V\leq2^9L^2k^{-3}$.

On the other hand, if $E<4V$, we have $kV-L\leq4V$, and hence $V\leq\frac{L}{k-4}$. As long as $k\geq8$, this implies $V\leq2L/k$, and we are done. Finally, for $k<8$, the trivial bound
\[
|P_k|\leq\binom L2/\binom k2\leq2L^2k^{-2}\leq2^9L^2k^{-3}.
\]
\end{proof}
